\chapter{HARDWARE IMPLEMENTATION OF THE DETECTOR}\label{ch:hw}

\section{Design and Interface}

The detector is one IP core inside the KV260 design (core version 3). It is reached through one AXI4-Lite slave at
100\,MHz with an 8-KB window: control and result registers, the feature registers, and the feature engine's register
file at address 0x1000 (Table~\ref{tab:core}).

\begin{table}[H]
\centering\small
\caption{Blocks of the SeizeLUT core}\label{tab:core}
\begin{tabularx}{\textwidth}{|P{34mm}|L|}
\hline
\multicolumn{1}{|c|}{\textbf{Block}} & \multicolumn{1}{c|}{\textbf{Function}}\\
\hline
Feature engine & Computes the 38 features from raw samples (Section~\ref{sec:fe-engine})\\
\hline
Classifier & 146 threshold comparators, 50 LUT6, popcount score, frozen and programmable operating points, 10-of-10
alarm; latency 3 clock cycles\\
\hline
Self-test (BIST) & Vector ROM and response check; 1,024 vectors in 1,028 cycles\\
\hline
Fault injection & Registers that flip any truth-table bit or force any comparator or LUT output to 0 or 1\\
\hline
Debug read-back & All 146 comparator bits and 50 LUT outputs of the last window\\
\hline
Demo I/O & PMOD outputs for the alarm LED, buzzer, score LEDs and a heartbeat; a debounced button input\\
\hline
\end{tabularx}
\end{table}

\section{Feature Engine in the Fabric}\label{sec:fe-engine}

A front end with dozens of filters would need many multipliers if built directly. Instead, a small microcoded
\emph{feature engine} computes the 38 features the classifier reads, one instruction at a time:

\begin{itemize}
\item one 40-bit adder and two 5-bit shifters; every filter coefficient is a sum of at most five signed powers of
  two, so a multiplication becomes a few shifted additions;
\item a register file of 930 words and a 2,323-instruction program in block RAM, at 3 clock cycles per instruction;
\item four events: a 256-Hz biopotential sample, a 25-Hz motion sample, the end of a second, and the last window.
\end{itemize}

After each window the core copies the engine's 38 features into the classifier and classifies them automatically. The
engine needs under one million cycles per second of signal, so it keeps up with the sensors at 1\,MHz, a hundredth of
the board's clock.

\section{Operation and Verification}

\begin{enumerate}
\item The Arm writes the raw samples of one second into the engine's register file and starts the events.
\item The engine runs its program and writes the 38 features; the core copies them into the classifier registers.
\item The comparators, the 50 LUTs and the popcount produce the score, the decision and the alarm in 3 cycles.
\item The Arm reads the result; the LED, buzzer and ESP32 wire follow the alarm bit directly.
\end{enumerate}

Every stage was verified against the software model before the board build: the engine program in an
instruction-level simulator, the engine RTL sample by sample and the whole core through the AXI bus in Verilator
(Section~\ref{sec:verif}).

\section{Built-In Self-Test}\label{sec:bist}

\textbf{Fault model.} Every single fault in the part of the hardware the network adds: each of the $50 \times 64 =
3{,}200$ truth-table bits flipped (an upset in FPGA configuration memory) and each comparator and LUT output stuck at 0
or 1 (292 + 100), 3,592 faults in total.

\textbf{Which faults matter.} Each fault was applied to all 31.4 million windows (8,714 hours) of the 96 training
patients. A fault is \emph{decision-relevant} if it changes at least one decision, and \emph{alarm-relevant} if it
changes at least one alarm.

\textbf{Test vectors.} 255 test vectors were chosen from training windows for coverage, so that together they
exercise every decision-relevant fault. The on-chip self-test runs 1,024 vectors in 1,028 clock cycles
(10.3\,\textmu s), so the device can test itself after every one-second window. The detection results are in
Section~\ref{sec:res-bist}.
